%% RESUMO EM LINGUA ESTRANGEIRA (elemento obrigatorio -- NBR 6028)
%% Traducao fiel do resumo em portugues. As keywords sao definidas em
%% config/dados.tex. Versao provisoria (dissertacao dummy).

\begin{abstract}
Audiovisual educational repositories maintained by public school systems
constitute digital assets of the State, whose public value depends less on their
existence than on the institutional capacity to convert them into effective
pedagogical use. This research develops and evaluates an artificial
intelligence based pedagogical recommendation artifact that classifies the 760
active video lessons of the ENEM preparatory course of Canal Educação, a program
of the Piauí State Department of Education, according to the 120 skills of the
reference matrix of the Brazilian National High School Examination (ENEM), in
order to support teacher planning without replacing pedagogical judgment.
Design Science Research is adopted as the method. The artifact converts
audiovisual content into text through automatic speech recognition, represents
lessons and curricular descriptors through embeddings of the multilingual
BGE-M3 model, used exclusively in inference, and ranks, for each lesson, the
closest skills by cosine similarity, without relying on usage history. The
technical evaluation compares the artifact method, a lexical method without
artificial intelligence (TF-IDF) and a generative model against the official
assignment of skills to exam items published by INEP. On the 178 items of the
2023 edition, the artifact outperforms the lexical method at the first ranking
position and within five positions, but reaches less than half of the agreement
obtained by the generative model. The results indicate the technical
feasibility of content-based recommendation in a cold-start context and reveal
limitations that condition the interpretation of the outputs. From a public
management standpoint, the classification of the repository becomes an input
for the governance of the asset itself, by exposing the curricular coverage
gaps that should guide the prioritization of new content production.
\end{abstract}
